% Appendix A.  Proofs of the population spectral geometry.  The order lem:A-decomp ->
% prop:coherence -> lem:gap -> lem:A-decomp(iv) is load-bearing: lem:gap uses only the
% inequality lambda_2 <= m*beta_0, and clause (iv) then uses lem:gap's interlacing
% bound to force equality.  Reordering makes the identification circular.
\section{Proof of the spectral geometry}\label{app:comp1}

\noindent Every later appeal to a bounded population entry is to \Cref{prop:cfbm}.
No range claim is made for $\hat S$: by \eqref{eq:sampling_estimator} its observed
off-diagonal entries equal $S_{ij}/p$, which exceed $1$ throughout the regime
$p=\Theta(\log m/m)$ of interest.

\noindent We first verify that the model of \Cref{def:cfbm} is realised at the root split.

\begin{proof}[Proof of \Cref{prop:cfbm}]
Let $i \in A$ and $j \in B$. Every path from $A$ to $B$ traverses the root, so
$\mathrm{MRCA}(i,j) = r$ and $\tau_{ij} = \tau_r$ for every cross-clan pair; this is purely
topological. By \eqref{eq:similarity_clock} all cross-clan entries then equal the single value
$\bze = \exp(2\,\mathrm{tr}(Q)\tau_r)$, which lies in $(0,1)$ because $\mathrm{tr}(Q)<0$ and
$\tau_r>0$.

Now let $i \ne j$ lie in the same clan, say $A$, with clan root $h_A$. Then $\mathrm{MRCA}(i,j)$ is
a descendant of $h_A$, so $\tau_{ij} \le \tau^{\max} \le \max(\tau_{h_A},\tau_{h_B}) < \tau_r$, the
last inequality strict because the two root branches have positive length. As $\mathrm{tr}(Q)<0$,
\eqref{eq:similarity_clock} is strictly decreasing in $\tau_{ij}$, so
$S_{ij} \ge \exp(2\,\mathrm{tr}(Q)\tau^{\max}) = \Sin^{\min} > \bze$. In the other direction,
$\mathrm{MRCA}(i,j)$ is a strict ancestor of both $i$ and $j$ for $i \ne j$, so $\tau_{ij}>0$ and
$S_{ij}<1$; taking the maximum over the finitely many pairs gives $\Smax<1$. Together with
$\tau_{ii}=0$, giving $S_{ii}=\exp(0)=1$, this is \eqref{eq:cfbm_floor} and the stated range. The
block degrees follow by summing the rows of $S$ within each block and adding the constant
cross-block contribution, $|B|\bze$ for a row of $A$ and $|A|\bze$ for a row of $B$.
\end{proof}

\noindent We next prove the spectral decomposition stated in \Cref{sec:bridging},
which the proof of \Cref{lem:gap} then uses.

\begin{proof}[Proof of \Cref{lem:A-decomp}(i)--(iii)]
For any $L=D-S$, $L\1=D\1-S\1=0$, so $0$ is an eigenvalue with
$v^{(1)}=\tfrac1{\sqrt m}\1_m$; no assumption on $S$ is needed. Every off-diagonal
entry of $S$ is at least $\bze>0$ by \Cref{prop:cfbm}, so the similarity graph is
connected and $\lambda_1=0$ is simple.

\emph{Step heights.} Assume $\vv$ has the stated form with $c_A,c_B>0$. The
constraints $\norm{\vv}_2=1$ and $\langle v^{(1)},\vv\rangle=0$ give $|A|c_A=|B|c_B$,
hence
\[
   |A|\Bigl(\tfrac{|B|}{|A|}c_B\Bigr)^2+|B|c_B^2=1
   \;\Longrightarrow\;\tfrac{|B|m}{|A|}c_B^2=1
   \;\Longrightarrow\; c_B^2=\tfrac{|A|}{|B|m},\quad c_A^2=\tfrac{|B|}{|A|m}.
\]

\emph{Eigenpair verification.} Write $\vv=(c_A\1_{|A|},\,-c_B\1_{|B|})^\top$ with
$|A|c_A=|B|c_B$. The only structural input is that the diagonal sub-degrees of
\Cref{prop:cfbm} already carry the row sums of their own block,
$D_A=\operatorname{diag}(S_A\1_{|A|})+|B|\bze I_{|A|}$, so that
\begin{equation}\label{eq:block_annihilation}
   (D_A-S_A)\1_{|A|}=|B|\bze\1_{|A|},\qquad(D_B-S_B)\1_{|B|}=|A|\bze\1_{|B|},
\end{equation}
\emph{identically, for arbitrary symmetric $S_A,S_B$}: no constancy of the diagonal
blocks is used, and the intra-clan floor of \eqref{eq:cfbm_floor} plays no role here.
Using the block form of $L$ with $J_{|A|\times|B|}\1_{|B|}=|B|\1_{|A|}$,
$J_{|B|\times|A|}\1_{|A|}=|A|\1_{|B|}$ and \eqref{eq:block_annihilation},
\[
   Lv^{(2)}=
   \begin{pmatrix}
     |B|\bze c_A\1_{|A|}+|B|\bze c_B\1_{|A|}\\[2pt]
     -|A|\bze c_A\1_{|B|}-|A|\bze c_B\1_{|B|}
   \end{pmatrix},
\]
which reduces to
$\bigl(|B|\bze(c_A+c_B)\1_{|A|},\,-|A|\bze(c_A+c_B)\1_{|B|}\bigr)^\top$. From
$|A|c_A=|B|c_B$, $|B|(c_A+c_B)=mc_A$ and $|A|(c_A+c_B)=mc_B$, so $Lv^{(2)}=m\bze\,\vv$: the pair
$(m\bze,\vv)$ is an eigenpair, which is claim~(ii).

\emph{The bound \eqref{eq:lambda2_upper}.} Since $\lambda_1=0$ is simple with
eigenvector $v^{(1)}$, the Courant--Fischer characterisation of the second
eigenvalue reads $\lambda_2=\min\{x^\top Lx:\ \norm{x}_2=1,\ x\perp v^{(1)}\}$. The
vector $\vv$ is admissible in that minimisation --- it is a unit vector and
$\langle v^{(1)},\vv\rangle=0$ by construction --- so
$\lambda_2\le\langle\vv,L\vv\rangle=m\bze\norm{\vv}_2^2=m\bze$.
\end{proof}

\noindent Nothing above forces $\lambda_2$ to \emph{equal} $m\bze$; that is an
ordering claim, and it is what \Cref{asm:margin} supplies in clause (iv), proved at
the end of this appendix after the interlacing bound it needs.

\begin{proof}[Proof of \Cref{prop:coherence}]
Substituting the Fiedler coordinates of \Cref{lem:A-decomp} into
\eqref{eq:coherence_definition},
\[
   \coh=\tfrac m2\max\Bigl\{\tfrac1m+c_A^2,\ \tfrac1m+c_B^2\Bigr\}
       =\tfrac m2\Bigl(\tfrac1m+\tfrac{|B|}{m|A|}\Bigr)
       =\tfrac12\Bigl(1+\tfrac{|B|}{|A|}\Bigr)=\tfrac{1+\eta}{2}. \qedhere
\]
\end{proof}

\begin{proof}[Proof of \Cref{lem:gap}]
The argument uses \Cref{asm:clock} and \Cref{prop:cfbm}, and only the
\emph{inequality} \eqref{eq:lambda2_upper}, $\lambda_2\le m\bze=m\Sout^{\max}$ ---
not clause (iv) of \Cref{lem:A-decomp}, whose proof appeals to the interlacing step
below.
For $\lambda_3$, standard algebraic-connectivity interlacing gives
$\lambda_3(L)\ge\min\bigl(\lambda_2(L_{C_1}),\lambda_2(L_{C_2})\bigr)$, and applying
the Fiedler-eigenvalue formula to each sub-clan Laplacian gives
$\lambda_2(L_{C_c})\ge n_c\,S_{\mathrm{in}}^{(c)}$. Since
$\min_c S_{\mathrm{in}}^{(c)}=\Sin^{\min}$, the weakest intra-clan similarity of
\eqref{eq:cfbm_floor},
\begin{equation}\label{eq:lambda3_lower}
   \lambda_3\ \ge\ \min\bigl(\lambda_2(L_{C_1}),\lambda_2(L_{C_2})\bigr)
   \ \ge\ \min(n_1,n_2)\,\Sin^{\min}=m_{\min}\Sin^{\min},
\end{equation}
so $\Dlam=\lambda_3-\lambda_2\ge m_{\min}\Sin^{\min}-m\Sout^{\max}$. Using
$\Sin^{\min}=\rho+\Sout^{\max}$, which is \eqref{eq:margin_under_clock}, and
$m=m_{\min}(1+\eta)$,
\[
   \Dlam\ge m_{\min}(\rho+\Sout^{\max})-(1+\eta)m_{\min}\Sout^{\max}
        =m_{\min}\rho-\eta\,m_{\min}\Sout^{\max}. \qedhere
\]
\end{proof}


\begin{proof}[Proof of \Cref{cor:tolerance}]
Positivity of the bound in \Cref{lem:gap} requires
$\eta<\Sin^{\min}/\Sout^{\max}$. The minimum intra-clan similarity is bounded by
the diameter, $\Sin^{\min}\ge e^{-r(\mathcal T)}$, and the maximum cross-clan
similarity by the subtree depth, $\Sout^{\max}\propto e^{-\mathcal O(h(\mathcal T))}$.
Substituting into $\eta<\Sin^{\min}/\Sout^{\max}$ gives \eqref{eq:eta_distance_bound};
the logarithmic scaling $r,h=\mathcal O(\log m)$ under Kingman yields
$\eta\le\mathcal O(e^{\mathcal O(\log m)})=\mathcal O(m^c)$.
\end{proof}

\noindent With \eqref{eq:lambda3_lower} in hand the inequality
\eqref{eq:lambda2_upper} can be closed.

\begin{proof}[Proof of \Cref{lem:A-decomp}(iv)]
Suppose $\lambda_2<m\bze$. Then $m\bze$, an eigenvalue of $L$ by
\Cref{lem:A-decomp}(ii) and equal to neither $\lambda_1=0$ nor $\lambda_2$, satisfies
$m\bze\ge\lambda_3$. The interlacing bound \eqref{eq:lambda3_lower} gives
$\lambda_3\ge m_{\min}\Sin^{\min}$, so this forces $m_{\min}\Sin^{\min}\le m\bze$.
But $\Sout^{\max}=\bze$ under \Cref{asm:clock} and $m=m_{\min}(1+\eta)$, so
\eqref{eq:margin_assumption} reads $\Sin^{\min}>(1+\eta)\bze$, i.e.\
$m_{\min}\Sin^{\min}>m\bze$ --- a contradiction. Hence $\lambda_2=m\bze$, and then
$\lambda_3\ge m_{\min}\Sin^{\min}>m\bze=\lambda_2$ shows $\lambda_2$ is simple, so
$\vv$ is the Fiedler vector up to sign.
\end{proof}
